\section{Protecting tentative offers when recommendations change}\label{app:protect-detail}
A housing service may already know who holds a tentative offer. Can it still use
a learned score to choose other applicants' lists while protecting those offers?
Here protection means keeping a placement; the unit or its rent can change.
Let $M_0$ be stable on the original display. Keep the visible options and
relative preferences of every member of $I$ fixed, and write $H$ for everyone
else. Listing priorities over incumbents are fixed. Let $Q$ be the units
assigned to $I$ in $M_0$, and write $i(j)=M_0^{-1}(j)$ for their holders.
For proposed lists of $H$, draw an edge $hj$ when $j\in Q$, $j$ is displayed
to $h$, and $j$ could prioritize $h$ over $i(j)$ under the information available
to the platform. If priorities are unknown, keep every displayed edge into $Q$.
Write $C$ for this bipartite graph and $\nu(C)$ for its maximum matching size.

\begin{theorem}[A prospective bound on displaced renters]\label{thm:threat}
In every stable outcome on the proposed display, at most $\nu(C)$ previously
matched members of $I$ become unmatched. The bound is attained for every
nonnegative integer value, even at quota one and with a common priority. More
strongly, every bipartite threat graph admits an instance attaining its matching number.
\end{theorem}
\begin{proof}
Write $M_1$ for the new matching. Start at a previously matched incumbent $i$
who is now unmatched. They prefer their old unit $j=M_0(i)$ to their new outcome.
By stability at $M_1$, $j$ has a new holder ranked above $i$. If the holder is
an outsider $h$, we reach a selected edge of $C$. Otherwise the holder $i'$
is an incumbent. The pair $i'j$ was visible before the switch, with $j$ ranking
$i'$ above $i$. Stability at $M_0$ forces $i'$ to have had a partner strictly
preferred to $j$. Continue at $i'$. These walks follow alternating edges in
the symmetric difference of two matchings, whose components are paths and
cycles. The start vertex has degree one, so its component is a path. The walk
can stop only at an outsider. Distinct starting losers give distinct terminal
listings and distinct terminal outsiders: each path is traced from its unmatched
incumbent end to its first outsider, and every noninitial step has a unique
predecessor. Thus the terminal edges form a matching in $C$. For equality take
disjoint copies of one incumbent whose only unit is $j$ and one higher-priority
outsider who switches from a private fallback to $j$.

For sharpness conditional on the graph, give each $j\in Q$ a dedicated incumbent
who accepts only $j$. Choose a maximum matching of $C$. In a common priority put
its matched outsiders first, then all other outsiders, then all incumbents.
Initially every outsider sees only private fallback listings; every incumbent
holds their sole unit. On the new display each outsider sees all its neighbors in
$C$. Matched outsiders rank their respective assigned neighbors first. They take
these distinct listings in priority order. Any unmatched outsider has only
matched neighbors, by maximality, and displaces no additional incumbent.
Private fillers ranked below all neighbors ensure both displays give each
outsider exactly $k$ slots for any $k\ge\max\{1,\Delta_H(C)\}$.
\end{proof}

This is a gross-loss guarantee: it counts an original recipient who loses
their placement even if someone else from the same group gains one. It needs
the observed baseline offers, a stable baseline under the specified display
rule, and unchanged incumbent lists, preferences and within-incumbent priorities.
The priority-blind version treats all higher-priority completions as possible. Computing its graph also needs the feasible choices of each outsider;
it uses neither post-change assignments nor outsiders' rankings. For a single
deployment without baseline offers, the bound cannot be evaluated as stated.
A display can have $\nu(C)>0$ and retain everyone, for example if a threatened
incumbent has an uncontested second choice.

\begin{corollary}[A retained offer as a rank floor]\label{cor:guard-rank}
If $C$ has no edges and the new market uses applicant-proposing DA,
every originally matched member of $I$ receives a unit they weakly
prefer to their original assignment.
\end{corollary}
\begin{proof}
Run DA with one proposal at a time. Suppose some incumbent is rejected or
displaced by their own original unit, and take the first such rejection.
Let $i$ be that incumbent and $j=M_0(i)$ the unit. Unit $j$ holds an
applicant ranked above $i$. An outsider in that position would give a
threat edge. An incumbent $i'$ in that position proposed to $j$. Because
$i'j$ was visible originally and $j$ ranks $i'$ above $i$, original stability
requires $i'$ to have held a unit strictly preferred to $j$. To have proposed
to $j$, $i'$ must already have been rejected by that original unit, contrary
to the choice of the first rejection. Every incumbent's original unit remains
visible. An incumbent ending at a worse option would have received such a
rejection.
\end{proof}
Applicant-proposing DA matters for the rank floor. For two incumbents with
opposite top choices, let each unit rank its other applicant first. The
applicant-optimal and unit-optimal matchings are both stable. Switching
between them can make both incumbents worse off with no outsiders at all.
The rank floor still permits a higher rent for a more-preferred unit.

A budget of one already permits arbitrarily many rank losses with no one
becoming unmatched. For $n$ incumbents, applicant $i$ ranks $s_i$ over
$s_{i+1}$ for $1\le i\le n$. One outsider switches from their private
fallback $z$ to $s_1$. The common priority puts the outsider first, then
incumbents in increasing order. Initially all incumbents get their first
choice; on rollout each gets their second. Only the outsider's edge to $s_1$
is a threat. Giving the outsider the same private, last-ranked filler on both
displays preserves a universal quota of two.

With known priorities, the first-rejection argument also locates possible rank losses when $C$ has edges. Start at the holders of threatened units. From baseline holder $a$ draw an arrow to holder $b$ if $a$ can see $M_0(b)$ and outranks $b$ there. Every rejected holder is reachable from a start: a rejection by an outsider starts a path, and a rejection by another fixed applicant extends a path from that applicant's earlier rejection. Induction over the rejection times proves the claim. It explains why regions of the market can stay protected even if the overall budget is positive.

\paragraph{Choosing a display with a loss budget.}
Let $q_h=\min\{k,d_h\}$, and let $w_{hj}$ be any score estimated before the
switch. On all feasible outsider edges use binary display indicators $x_{hj}$.
On the full set $C_*$ of possible threats use binary cover indicators
$a_h,b_j$ at their endpoints. Solve
\[
 \max \sum_{hj}w_{hj}x_{hj}
 \quad\text{subject to}\quad
 \sum_j x_{hj}=q_h,\qquad
 x_{hj}\le a_h+b_j\ (hj\in C_*),\qquad
 \sum_h a_h+\sum_{j\in Q}b_j\le B .
\]
By K\H{o}nig's bipartite matching theorem, a proposed display admits such
a cover exactly when its certificate $\nu(C)$ is at most $B$. At zero budget
each outsider receives their $q_h$ highest-scoring feasible options outside
the threat edges; feasibility requires and suffices that every row has $q_h$
such choices. If no certified slate exists, retaining the entire baseline
display preserves quotas and placements even when the baseline graph itself
has positive certificate. The following refinement can let other applicants switch.

\paragraph{Preserving every quota when a row has too few safe options.}
Start with $J=I$. Hold every $J$ row at its original ordered display and
protect the baseline units of all matched members of $J$. For each $h\notin J$
count its feasible nonthreat edges under this enlarged fixed set. If any
applicant cannot fill their quota, add all such applicants to $J$, retaining
their original displays. Repeat until each remaining applicant can fill their
quota, then display their top-scoring nonthreat options. The procedure stops
within $|H|$ nonempty rounds and preserves every quota. Theorem~\ref{thm:threat}
at the final $J$ gives zero losses among the initially matched members of $I$. Corollary~\ref{cor:guard-rank} also protects their original ranks under applicant-proposing DA.
For closure, fix in advance which listing priority comparisons are known; if more applicants are frozen, the information used for a given comparison stays fixed. The final $J$ is contained in any superset of $I$ that would admit complete
quota-respecting lists and a zero-threat certificate for all its members.
Indeed, given $J\subseteq J'$, a row deficient at $J$ must belong to any such
$J'$: as long as this row remains outside, enlarging $J$ only removes safe
edges. Induction proves the claim. The minimum freeze set is relative to this
specific certificate; another stable display can sometimes protect all renters
with fewer freezes. If every applicant is frozen, the original display returns.

\paragraph{The cost of a positive loss budget.}
Deciding whether a display can reach score $W$ subject to a certificate budget
$B$ is NP-complete even at quota one, with binary scores, at most three feasible
listings per outsider and a common strict priority. Reduce from Vertex Cover.
Make a unit for each graph vertex, held by a dedicated incumbent whose only
option is that unit. Introduce an outsider for each graph edge, with its two
endpoint units scoring one and a private fallback scoring zero. Initially each
outsider sees their fallback, and every incumbent holds their own unit. Give
every outsider higher priority than every incumbent. Set $W$ equal to the
number of edges. Every outsider must now nominate one endpoint. The selected
threat graph has outsider-degree one, so its matching number equals the number
of distinct nominated endpoints. These form a vertex cover. Conversely any
vertex cover permits each outsider to nominate one of its endpoints. Both
displays have one entry per applicant. Verifying the score, quotas and
matching number takes polynomial time. Thus finding the best score under
a variable loss budget can require combinatorial optimization, while checking
the certificate for a given display is efficient.
